\begin{table*}[t!]
\small  % You can also use \footnotesize, \scriptsize, or \tiny
\centering
\begin{tabularx}{\textwidth}{|X|}
\hline
\#\#\# Demo 1\\ Question: How to augument the datapipe by repeating it six times.\\ \\ Candidate Documents\\ API\_1: map(*args, **kwds):\textbackslash{}n    Apply the input function over each item from the source DataPipe (functional name: ``map``).\textbackslash{}n    The function can be any regular Python function or partial object.\\ API\_2: EndOnDiskCacheHolder(datapipe, mode='wb', filepath\_fn=None, *, same\_filepath\_fn=False, skip\_read=False):\textbackslash{}n    Indicates when the result of prior DataPipe will be saved local files specified\textbackslash{}n    by ``filepath\_fn`` (functional name: ``end\_caching``). Moreover, the result of source DataPipe\textbackslash{}n    is required to be a tuple of metadata and data, or a tuple of metadata and file handle.\\ API\_3: cycle(*args, **kwds):\textbackslash{}n    Cycles the specified input in perpetuity by default, or for the specified number of times (functional name: ``cycle``).\\ \\ Judgment: A\\ \\ By using above documents, you can generate the code snippet like 'datapipe.cycle(6)'.\\ As documents support clear, accurate, and engaging answers, query is answerable so your output should be A.\\ \\ \#\#\# Demo 2\\ Question: Read the URL using the HTTP protocol and process the csv file.\\ \\ Candidate Documents\\ API\_1: IoPathFileLister(*args, **kwds):\textbackslash{}n    Lists the contents of the directory at the provided ``root`` pathname or URL, and yields the full pathname or URL for each file within the directory.\\ API\_2: OnlineReader(*args, **kwds):\textbackslash{}n    Takes file URLs (can be HTTP URLs pointing to files or URLs to GDrive files), and yields tuples of file URL and IO stream.\\ API\_3: HttpReader(source\_datapipe: IterDataPipe{[}str{]}, timeout: Optional{[}float{]} = None):\textbackslash{}n    Takes file URLs (HTTP URLs pointing to files), and yields tuples of file URL and IO stream.\\ \\ Judgment: P\\ \\ You can read the URL with API\_3 HttpReader, but you cannot process the csv file with above documents.\\ You can partially answer to the query, so your output should be P.\\ \\ \#\#\# Demo 3\\ Question: Clone the source datapipe two times\\ \\ Candidate Documents\\ API\_1: unbatch(datapipe: IterDataPipe, unbatch\_level: int = 1):\textbackslash{}n    Undoes batching of data (functional name: ``unbatch``). In other words, it flattens the data up to the specified level within a batched DataPipe.\\ API\_2: IterToMapConverter(*args, **kwds):\textbackslash{}n    Lazily load data from ``IterDataPipe`` to construct a ``MapDataPipe`` with the key-value pair generated by ``key\_value\_fn`` (functional name: ``to\_map\_datapipe``).\\ API\_3: cycle(*args, **kwds):\textbackslash{}n    Cycles the specified input in perpetuity by default, or for the specified number of times (functional name: ``cycle``).\\ \\ Judgment: U\\ \\ You may need information of methods like 'clone()' to resolve the question.\\ But there's no relevant information about the methods. You are unanswerble to the query, so your output should be U. \\ \hline
\end{tabularx}
\caption{3-way-1shot example for TorchDataEval}
\label{tab:2shot_prompt_torchdataeval}
\end{table*}